%% kapa: exam=1 round=21 year=2010 period=2 form=A subject=appraisal_law subject_ko=부동산관계법규 seq=19 number=59
%% source: Q-Net 감정평가사 기출문제 (cst00309) artlSeq=5024181 file=extracted/1008/5024181/2010년 제21회 감정평가사 1차시험 기출문제/2교시.hwp
%% conversion: kordoc 4.24.51 -> tools/gen1.py (auto); review=proofread-gpt-sol; flags=none
\begin{question}{59}
부동산 가격공시 및 감정평가에 관한 법령상 표준지공시지가의 결정에 관한 설명으로 옳지 않은 것은?
\begin{choices}
  \item 국토해양부장관은 토지이용상황이나 주변환경 그 밖의 자연적․사회적 조건이 일반적으로 유사하다고 인정되는 일단의 토지 중에서 표준지를 선정하여야 한다.
  \item 국토해양부장관은 표준지에 대하여 매년 공시기준일 현재의 적정가격을 조사․평가한 후 중앙부동산평가위원회의 심의를 거쳐 이를 공시하여야 한다.
  \item 표준지의 적정가격을 조사․평가할 때에는 인근유사토지의 거래가격․임대료 및 당해 토지와 유사한 이용가치를 지닌다고 인정되는 토지의 조성에 필요한 비용추정액 등을 종합적으로 참작하여야 한다.
  \item 표준지의 적정가격을 조사․평가할 때에는 둘 이상의 감정평가업자에게 이를 의뢰하여, 평가한 금액의 산술평균치를 기준으로 한다.
  \item 감정평가업자가 행한 표준지의 조사․평가가 관계법령을 위반하여 평가되었다고 인정되는 경우에는 당해 감정평가업자가 아닌 다른 감정평가업자로 하여금 다시 조사․평가하도록 하여야 한다.
\end{choices}
\end{question}
\answer{5}
